%% kapa: exam=1 round=33 year=2022 period=2 form=A subject=appraisal_law subject_ko=감정평가관계법규 seq=36 number=36
%% source: Q-Net 감정평가사 기출문제 (cst00309) artlSeq=5209997 file=raw/1008/5209997/2022년도 제33회 감정평가사 1차시험 2교시(감정평가 관계법규, 회계학).pdf
%% conversion: kordoc 4.24.51 -> tools/gen1.py (auto); review=proofread-sonnet; flags=none
\begin{question}{36}
동산·채권 등의 담보에 관한 법령상 동산담보권에 관한 설명으로 옳지 않은 것은?
\begin{choices}
  \item 담보목적물의 훼손으로 인하여 담보권설정자가 받을 금전에 대하여 동산담보권을 행사하려면 그 지급 전에 압류하여야 한다.
  \item 담보권자가 담보목적물을 점유한 경우에는 피담보채권을 전부 변제받을 때까지 담보목적물을 유치할 수 있지만, 선순위권리자에게는 대항하지 못한다.
  \item 동산담보권을 그 담보할 채무의 최고액만을 정하고 채무의 확정을 장래에 보류하여 설정하는 경우 채무의 이자는 최고액 중에 포함되지 아니한다.
  \item 약정에 따른 동산담보권의 득실변경은 담보등기부에 등기를 하여야 그 효력이 생긴다.
  \item 동일한 동산에 관하여 담보등기부의 등기와 「민법」에 규정된 점유개정이 행하여진 경우에 그에 따른 권리 사이의 순위는 법률에 다른 규정이 없으면 그 선후에 따른다.
\end{choices}
\end{question}
\answer{3}
